\clearpage
\section{From token IDs to a trainable neural LM}\label{sec:neural}
\subsection*{Lookup and next-token alignment}
Let $E\in\R^{K_{\mathrm{in}}\times d}$ be the input table. For token ID $i$,
lookup returns $E[i]$. Equivalently, for a one-hot row vector $q_i$,
$q_iE=E[i]$. The one-hot construction explains the operation; an embedding
lookup avoids materializing the mostly zero tensor. IDs have integer dtype,
while the trainable weights and selected vectors are floating-point tensors.

For a batch of $B$ sequences with $T$ input positions, IDs have shape
$(B,T)$ and embeddings have shape $(B,T,d)$. The L03 example uses a table
with ten rows and width four. The table shape and the input shape are
different objects; neither is inferred from the numerical magnitude of an ID.

\fig{token-loss}{The neural bigram computation keeps one input token per
prediction. Lookup changes the representation; it does not extend the context.}

For the sequence $[1,2,3,4,5]$, use inputs $[1,2,3,4]$ and targets
$[2,3,4,5]$. A second row $[5,6,7,8,9]$ gives inputs $[5,6,7,8]$ and
targets $[6,7,8,9]$. There are eight predictions, with last targets 5 and 9.
For a vocabulary of size ten, logits have shape $(2,4,10)$.

An untied neural bigram computes
\[
 e_t=E[x_t],\qquad z_t=W_oe_t+b_o,\qquad
 p(x_{t+1}\mid x_t)=\softmax(z_t),
\]
where $W_o\in\R^{K\times d}$. In the general case, $K_{\mathrm{in}}$ and
$K$ need not be equal: an input-only boundary symbol may occupy a lookup
row without being an allowed output. L03's ten-symbol toy uses matching sizes.

Flatten logits to $(BT,K)$ and targets to $(BT)$ in the same order before
computing cross-entropy. Sliced tensors may be non-contiguous; a reshape
can preserve their logical order while making a copy if needed. Mixing batch
and time orders would silently train on the wrong labels.

\takeaway{Write down shape, dtype, device, and meaning at every step.
Correct dimensions alone do not prove correct target alignment.}
\selfcheck{For $B=2$, $T=4$, $d=4$, and $K=10$, how many logits exist?
There are 80; only eight are the logits of the observed target classes.}

\nextpage{Stable probabilities and a connected computation graph}
For one target $y$ and logits $z\in\R^K$, the loss is
\[
 \ell(z,y)=-z_y+\log\sum_{j=1}^{K}e^{z_j}.
\]
Let $a=\max_jz_j$. Subtracting $a$ inside the exponentials gives the stable
log-sum-exp identity
\[
 \ell(z,y)=-(z_y-a)+\log\sum_j e^{z_j-a}.
\]
All exponent arguments are nonpositive, avoiding overflow from a large
common offset. For $z=(1000,1001)$ and the second class as target,
$\ell=\log(1+e^{-1})\approx0.3133$ rather than an expression involving
overflowing $e^{1000}$ and $e^{1001}$.

\fig{loss-graph}{Autograd follows the operations connecting parameters to the
loss. Detaching the loss breaks that connection even when its printed value is unchanged.}

The gradient with respect to logits has a simple form:
\[
 \frac{\partial\ell}{\partial z_j}=p_j-\mathbf1[j=y].
\]
It increases pressure on the target score and decreases pressure on the
others. Summing these derivatives gives zero because probabilities sum to
one. Adding a common constant to all logits leaves probabilities unchanged.

With all logits zero, $p_j=1/K$ and the loss is exactly $\log K$.
For $K=10$, this is about 2.3026 nats. Random logits generally do not yield
that exact value; use an intentional zero-logit setup for the check.

\begin{lstlisting}
# logits: (B, T, K); targets: (B, T), integer IDs
loss = F.cross_entropy(
    logits.reshape(-1, logits.shape[-1]),
    targets.reshape(-1),
)
loss.backward()
\end{lstlisting}
Pass raw logits to this loss, which performs the stable normalization
internally. Calling \texttt{softmax} first changes the function being
optimized. Calling \texttt{loss.item()} returns a Python number for
logging; it cannot carry the computation graph. Calling \texttt{detach()}
similarly removes the path needed for this backward call.

\selfcheck{If training loss is numerically finite but every parameter gradient
is absent, what would you inspect first? Trace the computation graph and
the gradient-recording context before changing the learning rate.}

\nextpage{A feedforward model can use more than one token}
Consider ``red key'' and ``blue key.'' If the next word should be
\texttt{opens} in the first context and \texttt{closes} in the second,
a bigram model cannot separate them: its input is \texttt{key} in both cases.
An embedding with more dimensions still receives the same input ID.

\fig{window-model}{An ordered two-token context distinguishes the two prefixes.
The hidden layer mixes the concatenated embeddings before producing vocabulary logits.}

A neural probabilistic language model uses a fixed window of $m$ tokens.
With column-vector conventions,
\[
 u=\operatorname{concat}(E[x_{t-m}],\ldots,E[x_{t-1}]),\quad
 h=\tanh(W_hu+b_h),\quad z=W_oh+b_o.
\]
Here $u\in\R^{md}$, $W_h\in\R^{H\times md}$, and
$W_o\in\R^{K\times H}$. Concatenation preserves slot order. The nonlinear
activation allows behavior beyond a single affine map of the concatenated
inputs; composing affine layers without an activation would collapse to
one affine map of $u$.

The L04 corpus contains two separate sequences:
\begin{center}
\texttt{BOS red key opens EOS}\qquad\texttt{BOS blue key closes EOS}.
\end{center}
For $m=2$, each supplies three windows. For the first sentence they are
$(\BOS,\texttt{red})\to\texttt{key}$,
$(\texttt{red},\texttt{key})\to\texttt{opens}$, and
$(\texttt{key},\texttt{opens})\to\EOS$. The second sentence supplies
the analogous blue/closes examples. No window crosses a document boundary.

\begin{tabularx}{\textwidth}{@{}l Y@{}}\toprule
Stage & Shape for a two-example probe, $m=2,d=4,H=8,K=7$\\\midrule
Context IDs & $(2,2)$\\
Lookup / concatenate & $(2,2,4)$ / $(2,8)$\\
Hidden / logits & $(2,8)$ / $(2,7)$\\\bottomrule
\end{tabularx}

This toy allocates all seven listed symbols as both input and output rows,
including \BOS{}, although \BOS{} never appears as a target. That convention
explains its $K=7$ baseline. A model with \BOS{} excluded from the output
would need a different output mapping and denominator.

\takeaway{Shared learned representations improve generalization across
contexts, but a fixed window still has a hard history limit. See \readingref{R04}.}
\selfcheck{Which parameter matrix grows when the context window doubles while
$d,H,K$ stay fixed? The input-to-hidden matrix $W_h$ doubles its column count.}

\nextpage{Training, tiny-batch fitting, and evaluation}
Gradient descent adjusts parameters to reduce a chosen objective. A
training step must compute a connected loss, differentiate it, and apply
an update. Gradients normally accumulate, so clear old gradients when
each step is intended to use only the current batch.

\begin{lstlisting}
model.train()
for step in range(200):
    optimizer.zero_grad(set_to_none=True)
    logits = model(train_contexts)
    loss = F.cross_entropy(logits, train_targets)
    loss.backward()
    optimizer.step()
\end{lstlisting}
This loop assumes the model, optimizer, and sentence-local training windows
from the course notebook have already been constructed. The call order is
the transferable idea; the batch and objective determine what it learns.

\fig{training-curves}{Saved L03 and L04 teaching runs, redrawn from separate
assets. Left: neural bigram, $K=10,d=4$, eight pairs, seed 0. Right: two-token
NPLM, $K=7,d=4,H=8$, six windows, seed 7. Both use 200 full-batch SGD updates
with learning rate $0.5$ on CPU. These curves are not a controlled model comparison.}

L03's mean loss reaches about 0.0187 nats, while L04's reaches about
0.0070. Their vocabularies, batches, architectures, and seeds differ.
Both examples establish that these implementations can fit their small
compatible training sets. Neither curve measures generalization to new documents.

An incompatible tiny batch can expose a model limitation. If a bigram sees
the same token with two equally frequent targets, its best conditional
distribution splits probability between them. Its optimum for those two
events has loss $\log2$, not zero. A failed fit is therefore not automatically
an optimizer bug: first check whether the model's inputs distinguish the labels.

For evaluation, switch behavior with \texttt{model.eval()} and separately
disable gradient recording with \texttt{torch.no\_grad()}. The first affects
modules such as dropout; the second controls graph construction. Neither
performs an optimizer step. The L03 unseen-ID probe tests coverage of selected
input rows; it is a diagnostic, not a representative held-out corpus.

\selfcheck{A training loss falls while held-out loss rises. Which observation
supports fitting, and which warns about generalization? Report both with their splits.}

\nextpage{Weight tying and a memory ledger}
If the same $K$ symbols occupy input and output rows and the state width is
$d$, the output matrix can share the input table: $W_o=E$. Then
$z=Eh+b_o$. This saves a table, but changes the gradient paths. Compatible
row meanings and widths are essential; equal tensor shapes alone are insufficient.

\fig{weight-tying}{Untied tables have independent parameters. Tying makes one
parameter receive the sum of its input-lookup and output-readout gradients.}

For an untied lookup, a row unused as input has no gradient through that
lookup. With tying, it may receive a gradient through output normalization
even if its token never appeared in the input batch. In a finite-logit
softmax, every output class generally contributes. Do not count the same
shared parameter twice when reporting model size.

\fig{memory-ledger}{Payload for one $32{,}000\times512$ table. Each fp32 tensor
of this shape uses 62.5 MiB; two Adam moment tensors account for 125 MiB.
Activations and temporary allocations require additional memory.}

The table has $P=32{,}000\cdot512=16{,}384{,}000$ parameters.
At two bytes per parameter its bf16 payload is 31.25 MiB. If parameters,
gradients, and two Adam moment tensors all use fp32, their combined payload is
$4P+4P+8P=16P$ bytes, or 250 MiB. This is an explicit storage model;
mixed-precision implementations can store different copies and dtypes.

Output logits add a different cost. Materializing a $(B,T,K)$ tensor needs
$BTK$ elements. A dense output projection from width $d$ costs approximately
$2BTdK$ forward FLOPs when one multiply-add counts as two. In the L03 toy,
$B=2,T=4,d=4,K=10$ gives 80 logits and about 640 projection FLOPs.
This excludes softmax, backward computation, and the optimizer update.

Allocated embedding rows, tokenizer entry counts, and maximum token ID are
not always identical. Use the actual matrix shape for storage, and check
that every emitted ID falls within its valid row range.

\takeaway{Separate unique parameters, tensor payload, intermediate tensors,
and peak device memory. A parameter count alone does not specify a training budget.}
\selfcheck{If weights are tied, do output logits disappear? No. Sharing removes
one parameter table, while the per-position prediction computation remains.}
